\documentclass[11pt,a4paper]{article}
\usepackage[T1]{fontenc}
\usepackage[utf8]{inputenc}
\usepackage{libertinus}
\usepackage{microtype}
\usepackage[margin=26mm,headheight=15pt]{geometry}
\usepackage{amsmath,amssymb,amsthm,mathtools}
\usepackage{aliascnt}
\usepackage{booktabs,array,tabularx,longtable}
\usepackage{enumitem}
\usepackage{xcolor}
\usepackage{xurl}
\usepackage{hyperref}
\usepackage[nameinlink,noabbrev,capitalise]{cleveref}
\usepackage{fancyhdr}
\usepackage{listings}
\definecolor{ink}{HTML}{17324D}
\definecolor{accent}{HTML}{176B70}
\hypersetup{colorlinks=true,linkcolor=ink,citecolor=accent,urlcolor=accent,
 pdftitle={Sharp Stability Strata for Finite Fabius--Rvachev Deconvolution},
 pdfauthor={Research report prepared for Vladimir Reshetnikov},
 pdfsubject={Optimal inverse stability, convolution scales, cumulants, polynomial roots}}
\pagestyle{fancy}
\fancyhf{}
\fancyhead[L]{\small Finite Fabius--Rvachev deconvolution}
\fancyhead[R]{\small\thepage}
\renewcommand{\headrulewidth}{0.3pt}
\setlength{\parskip}{0.35em}
\setlength{\emergencystretch}{2em}
\setlist{itemsep=3pt,topsep=5pt}
\allowdisplaybreaks[2]
\numberwithin{equation}{section}
\setcounter{tocdepth}{1}
\newtheorem{theorem}{Theorem}[section]
\newaliascnt{lemma}{theorem}
\newtheorem{lemma}[lemma]{Lemma}
\aliascntresetthe{lemma}
\crefname{lemma}{Lemma}{Lemmas}
\newaliascnt{proposition}{theorem}
\newtheorem{proposition}[proposition]{Proposition}
\aliascntresetthe{proposition}
\crefname{proposition}{Proposition}{Propositions}
\newaliascnt{corollary}{theorem}
\newtheorem{corollary}[corollary]{Corollary}
\aliascntresetthe{corollary}
\crefname{corollary}{Corollary}{Corollarys}
\theoremstyle{definition}
\newaliascnt{definition}{theorem}
\newtheorem{definition}[definition]{Definition}
\aliascntresetthe{definition}
\crefname{definition}{Definition}{Definitions}
\newaliascnt{example}{theorem}
\newtheorem{example}[example]{Example}
\aliascntresetthe{example}
\crefname{example}{Example}{Examples}
\newaliascnt{problem}{theorem}
\newtheorem{problem}[problem]{Research question}
\aliascntresetthe{problem}
\crefname{problem}{Research question}{Research questions}
\theoremstyle{remark}
\newaliascnt{remark}{theorem}
\newtheorem{remark}[remark]{Remark}
\aliascntresetthe{remark}
\crefname{remark}{Remark}{Remarks}
\newcommand{\R}{\mathbb R}
\newcommand{\C}{\mathbb C}
\newcommand{\N}{\mathbb N}
\newcommand{\E}{\mathbb E}
\newcommand{\Law}{\mathcal L}
\newcommand{\TV}{d_{\mathrm{TV}}}
\newcommand{\up}{\operatorname{up}}
\newcommand{\sinc}{\operatorname{sinc}}
\newcommand{\supp}{\operatorname{supp}}
\newcommand{\dd}{\,\mathrm d}
\newcommand{\norm}[1]{\left\lVert#1\right\rVert}
\newcommand{\abs}[1]{\left|#1\right|}
\newcommand{\coef}{\operatorname{coef}}
\newcommand{\dist}{\operatorname{dist}}
\newcommand{\rank}{\operatorname{rank}}
\newcommand{\kap}{\kappa}
\newcommand{\sr}{\mathcal S_r}
\newcommand{\stirling}[2]{\genfrac{\{}{\}}{0pt}{}{#1}{#2}}
\newcommand{\err}{\mathcal E_r}
\newcommand{\base}{a^{\circ}}
\title{\textbf{Sharp Stability Strata for Finite\\Fabius--Rvachev Deconvolution}\\[0.6em]
\large Colliding scales, vanishing factors, and optimal recovery exponents}
\author{Research report prepared for Vladimir Reshetnikov}
\date{28 September 2026}

\begin{document}
\maketitle
\begin{abstract}
Consider an independent sum of a known smooth compactly supported random variable
and at most $r$ unknown centered uniform factors.  The Rvachev up distribution is
an important special background, connecting the problem directly to the infinite
sinc products and exact arithmetic developed in the ProveIt repository.  We prove
a complete local classification of the optimal inverse H\"older exponent in total
variation distance.  At a configuration with positive squared-scale multiplicities
$m_1,\ldots,m_s$ and $m_0$ zero slots, the optimal exponent for two varying nearby
configurations is
\[
 \frac1M,\qquad M=\max\bigl(\{m_1,\ldots,m_s\}\cup
             \{2m_0:m_0>0\}\bigr).
\]
The estimate uses only the first $r$ even cumulants for its upper bound.  Sharpness
is proved by explicit Chebyshev level-set constructions whose low-order power
sums cancel.  A vanishing $m$-factor cluster has a nonzero, explicitly identified
leading density difference of order $t^{2m}$.

A distinct theorem concerns comparison with the fixed base configuration: its
optimal exponent is $1$ at a simple positive configuration and $1/2$ at every
singular configuration, regardless of multiplicity.  Thus pointwise stability
at a singular point is substantially stronger than uniform pairwise stability
near it.  We also give exact moment reconstruction, deterministic optimal-recovery
consequences, noise-safe feasible fitting, an elementary sampling upper bound,
and a formalization roadmap.  Classical cumulant identities, Newton identities,
and polynomial-root perturbation are distinguished from the classification
proved here.  The package contains exact algebraic checks and high-precision
Fourier diagnostics; it does not claim a Lean verification, a statistical
minimax lower bound, or worldwide priority.
\end{abstract}
\newpage
\tableofcontents
\newpage

\section{The research question and the contribution boundary}
\label{sec:scope}

\subsection{From exact identifiability to quantitative recovery}
A product of sinc functions remembers its convolution factors.  Exact recovery,
however, is not the same question as stable recovery from inexact data.  If two
unknown factors have almost the same scale, or if several factors are becoming
invisible because their scales approach zero, arbitrarily small changes in the
observed law can conceal much larger changes in its factorization.  The purpose
of this article is to determine the exact power law of that loss of information
when the number of unknown factors is bounded.

The ProveIt Fabius development supplies a particularly natural setting.  Its
module \path{GeneralizedRvachevIdentifiability.lean} recovers admissible dyadic
exponent sequences from analytic orders at dyadic points and establishes
injectivity of the corresponding entire products \cite{RepoIdent}.  That theorem
uses zero multiplicities, not a noise metric on probability laws.  The present
problem allows finitely many arbitrary nonnegative scales, places them above a
known smooth background, and asks for optimal stability in total variation.
These are different hypotheses and different conclusions; the result below is
not advertised as a replacement for the repository's dyadic theorem.

The finite box convolutions and the infinite up product also occur in the
repository's holonomic study \cite{RepoHolonomic}.  The smooth compactly supported
function and its probabilistic interpretation have a much older literature,
including Arias de Reyna's account \cite{Arias}.  We will establish directly all
properties of the up background required by our proofs.

\subsection{What is classical and what is proved here}
The reduction from cumulants to power sums, and then to a polynomial by Newton's
identities, is classical algebra.  Collision singularities in moment inversion,
Vieta maps, and Prony systems are extensively studied; Batenkov and Yomdin give
a relevant geometric treatment \cite{Prony}.  The $n$th-root scale in degree-$n$
polynomial-root perturbation is also classical; see Sch\"atzle \cite{Schaetzle}.
Our root-matching lemma is deliberately proved below rather than imported with
an unspecified constant or an ambiguous treatment of multiplicities.

The principal result developed here is the combined classification for the
\emph{full output-law metric}: positive collisions contribute their multiplicity,
zero clusters contribute twice their multiplicity, and the largest contribution
controls the optimal local pairwise exponent.  The fixed-base exponent has a
different classification.  Explicit smooth-density sharpness constructions and
an exact vanishing-cluster leading term make both statements two-sided results,
not merely condition-number heuristics.

This is a focused research contribution, not a claim to have settled all
infinite-product inverse problems or a named longstanding conjecture.  The source
comparison comprised selected repository documents and targeted code searches,
not an exhaustive audit of every report.  No claim of worldwide novelty is made.
The mathematical claims labeled theorems have proofs in this article; those
proofs have not been independently refereed or checked by a proof assistant.

\subsection{A precise repository snapshot}
The pinned repository tree used for the principal source comparison was
\begin{center}
\small\texttt{0e0368b216b302d5822890b48bd19f2337af4d04}.
\end{center}
The identifiability module and the holonomic report were read at that tree.
Some search-index results referred to a different indexed tree, so search
hits were used for discovery rather than as evidence that all files had the
same version.  The accompanying \texttt{SOURCE\_AUDIT.md} records the scope and
bibliographic checks.  No repository files were changed.

\section{Model, metrics, and the principal theorems}
\label{sec:main}

\subsection{The fixed-background finite-factor model}
Throughout, $r\ge1$ is a fixed integer.  Let $B$ be a known random variable,
independent of independent variables $U_1,\ldots,U_r$, each uniform on $[-1,1]$.
Assume initially that $B$ has a probability density
\begin{equation}
 g\in C_c^\infty(\R),\qquad \supp g\subseteq[-L,L].
 \label{eq:background}
\end{equation}
Symmetry of $g$ is not required.  Define the closed ordered parameter cone
\begin{equation}
 \sr=\{a\in\R^r:0\le a_1\le\cdots\le a_r\},\qquad
 X_a=B+\sum_{i=1}^r a_iU_i,
 \quad \mu_a=\Law(X_a).
 \label{eq:model}
\end{equation}
Its density is
\begin{equation}
 f_a(x)=\E g\left(x-\sum_i a_iU_i\right).
 \label{eq:density}
\end{equation}
A zero coordinate contributes the deterministic value zero.  Thus $r$ is a
\emph{capacity}: the model contains all sums with at most $r$ positive uniform
factors, padded to length $r$.  There is no attempt to identify how many additional
formal zero factors someone may have written outside this fixed representation.
Ordering removes permutation ambiguity.  Restricting to nonnegative amplitudes
removes the sign ambiguity of a centered uniform distribution.

For $a,b\in\sr$, put
\begin{equation}
 d(a,b)=\max_i|a_i-b_i|,
 \qquad \TV(\mu_a,\mu_b)=\frac12\norm{f_a-f_b}_{L^1}.
 \label{eq:metrics}
\end{equation}
The first metric is the bottleneck matching distance on the two multisets.  For
real ordered lists, matching in increasing order minimizes the largest matching
distance: whenever two matched pairs cross, uncrossing them cannot increase that
largest distance.  Repeating this operation proves the assertion.

We use squared amplitudes and their power sums:
\begin{equation}
 p_i=a_i^2,\qquad s_k(a)=\sum_{i=1}^r p_i^k,\qquad
 \err(a,b)=\max_{1\le k\le r}|s_k(a)-s_k(b)|.
 \label{eq:power-metric}
\end{equation}
The normalization of this finite-data error assumes a fixed unit of scale;
rescaling the physical variables changes the constants in our estimates.

\subsection{Collision multiplicities}
Fix $\base\in\sr$.  Its squared coordinates have distinct positive values
$c_1<\cdots<c_s$, with respective multiplicities $m_1,\ldots,m_s$.
Let $m_0$ be the number of zero coordinates.  Thus
$m_0+\sum_{j=1}^s m_j=r$.  Define
\begin{equation}
 M(\base)=\max\left(\{m_j:1\le j\le s\}\cup
                  \{2m_0:m_0>0\}\right).
 \label{eq:M}
\end{equation}
The set in this maximum is nonempty, including when all coordinates vanish.

\begin{definition}[Pairwise versus anchored exponents]
A pairwise exponent $\alpha>0$ is valid near $\base$ if some relative neighborhood
$\mathcal U\subseteq\sr$ and constant $C$ satisfy
\[
 d(a,b)\le C\TV(\mu_a,\mu_b)^\alpha\quad(a,b\in\mathcal U).
\]
An anchored exponent is valid at $\base$ if the analogous estimate is required
only with $b=\base$.  ``Optimal'' means the largest valid exponent.
\end{definition}

\begin{theorem}[Complete local pairwise classification]
\label{thm:pairwise}
Under \eqref{eq:background}, the optimal pairwise exponent near $\base$ is
$1/M(\base)$.  More precisely, there are a relative neighborhood $\mathcal U$ and
positive constants $C_1,C_2$ such that
\begin{equation}
 d(a,b)\le C_1\err(a,b)^{1/M(\base)}
         \le C_2\TV(\mu_a,\mu_b)^{1/M(\base)}
 \quad(a,b\in\mathcal U).
 \label{eq:main-bound}
\end{equation}
In every relative neighborhood of $\base$, a bound for $d(a,b)$ by a constant
times either $\err(a,b)^\alpha$ or $\TV(\mu_a,\mu_b)^\alpha$ fails for every
$\alpha>1/M(\base)$.
\end{theorem}

\begin{theorem}[Complete anchored classification]
\label{thm:anchored}
The optimal anchored exponent is $1$ when all coordinates of $\base$ are positive
and pairwise distinct.  In every other case it is $1/2$.  In particular, if there
is a repeated positive scale or at least one zero slot, then
\begin{equation}
 d(a,\base)\le C\TV(\mu_a,\mu_{\base})^{1/2}
 \label{eq:anchored}
\end{equation}
for all sufficiently nearby $a$, and no larger anchored exponent is valid.
\end{theorem}

\begin{theorem}[Uniform bounds on bounded parameter sets]
\label{thm:global}
For every $A>0$, there is a constant $C$ such that
\begin{equation}
 d(a,b)\le C\TV(\mu_a,\mu_b)^{1/(2r)}
 \quad\text{whenever }0\le a_1\le\cdots\le a_r\le A.
 \label{eq:global}
\end{equation}
The exponent $1/(2r)$ is optimal on this set.  On the smaller set
$a_*\le a_1\le\cdots\le a_r\le A$, where $0<a_*<A$, the optimal uniform exponent
is $1/r$.  If in addition the squared coordinates have a fixed positive minimum
separation, the inverse is uniformly Lipschitz on that separated set.
\end{theorem}

\subsection{Interpretation and an example table}
The exponent measures a loss of resolution.  With error level $\varepsilon$ in
the observed law, a cluster of $m$ positive factors can force amplitude uncertainty
of order $\varepsilon^{1/m}$.  A cluster of $m$ factors disappearing at zero can
force uncertainty of order $\varepsilon^{1/(2m)}$.  These orders are attained by
explicit pairs of probability densities, not by arbitrary coefficient vectors
that might fail to correspond to real nonnegative scales.

\begin{center}
\begin{tabular}{@{}p{0.48\textwidth}ccc@{}}
\toprule
Configuration with capacity $r=3$ & $M$ & Pairwise & Anchored\\
\midrule
Three distinct positive scales & 1 & $1$ & $1$\\
A repeated positive pair and a different positive scale & 2 & $1/2$ & $1/2$\\
Three equal positive scales & 3 & $1/3$ & $1/2$\\
One zero and two distinct positive scales & 2 & $1/2$ & $1/2$\\
Two zeros and one positive scale & 4 & $1/4$ & $1/2$\\
Three zeros & 6 & $1/6$ & $1/2$\\
\bottomrule
\end{tabular}
\end{center}
For a mixed example with three equal positive scales and two zero slots,
$M=\max(3,4)=4$.  Vanishing factors, rather than the positive collision, control
the pairwise exponent.  Constants depend on the background, the base
configuration, the neighborhood, and the capacity; the theorems do not suppress
this dependence.

\section{The Fabius--Rvachev background and exact cumulants}
\label{sec:fabius}

\subsection{A self-contained construction of the required background}
Let $V_1,V_2,\ldots$ be independent uniform variables on $[-1,1]$ and define
\begin{equation}
 B_{\up}=\sum_{j=1}^{\infty}2^{-j}V_j.
 \label{eq:up-series}
\end{equation}
The series converges absolutely on every outcome and lies in $[-1,1]$.  We use
the radian convention $\sinc t=\sin(t)/t$, continued by $\sinc0=1$.
The characteristic function is
\begin{equation}
 \phi_{\up}(t)=\E e^{itB_{\up}}
       =\prod_{j=1}^{\infty}\sinc(2^{-j}t).
 \label{eq:up-product}
\end{equation}
Independence gives the finite products, and bounded convergence gives their limit.
For fixed $N$ and nonzero real $t$,
\begin{equation}
 |\phi_{\up}(t)|\le
 \prod_{j=1}^N\min(1,2^j/|t|)
 \le 2^{N(N+1)/2}|t|^{-N}.
 \label{eq:rapid-decay}
\end{equation}
Consequently $t^k\phi_{\up}(t)$ is integrable for every fixed $k$.  Fourier inversion
and differentiation under the integral give a $C^\infty$ density $\up$.  The
support restriction of the random series gives $\supp\up\subseteq[-1,1]$.
It is therefore a $C_c^\infty$ probability density, exactly the regularity needed
in \eqref{eq:background}.  This is the normalization of the up distribution used
in the cited repository report \cite{RepoHolonomic}; the construction above also
fixes it independently of notation in any other source.

In the convention $\widehat g(\xi)=\int e^{-2\pi i\xi x}g(x)\dd x$, the same
formula is $\widehat{\up}(\xi)=\prod_{h=0}^\infty\sinc(\pi\xi/2^h)$.
The shift of index and the factor $2\pi$ account for the apparent difference
between the two product displays.

\subsection{Uniform cumulants are nonzero in every even order}
For a compactly supported real random variable $Y$, define its cumulants by
\[
 \log\E e^{zY}=\sum_{n\ge1}\kap_n(Y)\frac{z^n}{n!}
\]
in a neighborhood of zero.  The logarithm is well defined there because the
moment generating function has value one at zero.  Independence adds cumulants,
and scaling by $a$ multiplies the $n$th cumulant by $a^n$.

\begin{lemma}[Uniform cumulants]
\label{lem:uniform}
If $U$ is uniform on $[-1,1]$, then its odd cumulants vanish and
\begin{equation}
 \gamma_{2k}:=\kap_{2k}(U)
 =\frac{2^{2k}B_{2k}}{2k}
 =(-1)^{k+1}\frac{(2k)!\,\zeta(2k)}{k\pi^{2k}}\ne0.
 \label{eq:uniform-cumulants}
\end{equation}
In particular, $\gamma_2=1/3$, $\gamma_4=-2/15$, and $\gamma_6=16/63$.
\end{lemma}
\begin{proof}
Direct integration gives $\E e^{zU}=\sinh z/z$.  Its logarithmic derivative is
$\coth z-1/z$.  The Bernoulli generating identity
$w/(e^w-1)=\sum_{n\ge0}B_nw^n/n!$ gives
\[
 \coth z-\frac1z
 =\sum_{k\ge1}\frac{2^{2k}B_{2k}}{(2k)!}z^{2k-1}.
\]
Integration with zero constant term and comparison with the cumulant expansion
prove the first expression.  Alternatively, the classical product
$\sinh z/z=\prod_{n\ge1}(1+z^2/(\pi^2n^2))$ gives
\[
 \log\frac{\sinh z}{z}
 =\sum_{k\ge1}\frac{(-1)^{k+1}\zeta(2k)}{k\pi^{2k}}z^{2k}
 \quad(|z|<\pi),
\]
by absolutely convergent logarithmic expansion.  This proves the second
expression and its nonvanishing, since $\zeta(2k)>0$.
\end{proof}

\begin{proposition}[Power-sum decoder]
\label{prop:decoder}
In the model \eqref{eq:model},
\begin{equation}
 s_k(a)=\frac{\kap_{2k}(\mu_a)-\kap_{2k}(B)}{\gamma_{2k}}.
 \label{eq:decoder}
\end{equation}
For the up background,
\begin{equation}
 \kap_{2k}(B_{\up})=\frac{\gamma_{2k}}{4^k-1},\qquad
 s_k(a)=\frac{\kap_{2k}(\mu_a)}{\gamma_{2k}}-\frac1{4^k-1}.
 \label{eq:up-decoder}
\end{equation}
\end{proposition}
\begin{proof}
Cumulant additivity and scaling give \eqref{eq:decoder}.  For the finite prefix
of \eqref{eq:up-series}, the $2k$th cumulant is
$\gamma_{2k}\sum_{j=1}^N2^{-2kj}$.  Its moments converge to those of $B_{\up}$
by uniform boundedness and almost sure convergence.  Cumulants are polynomials
in finitely many moments, so they also converge.  Summing the geometric series
gives \eqref{eq:up-decoder}.
\end{proof}
In particular, $\operatorname{Var}(B_{\up})=1/9$.  This normalization check is
useful: confusing the up variance with the variance $1/3$ of a single uniform
factor would give an incorrect decoder at its first step.

\section{Exact reconstruction and quantitative algebraic reduction}
\label{sec:algebra}

\subsection{Newton reconstruction}
For $p_i=a_i^2$, write
\begin{equation}
 P_a(z)=\prod_{i=1}^r(z-p_i)
       =z^r-e_1z^{r-1}+e_2z^{r-2}-\cdots+(-1)^re_r.
 \label{eq:poly}
\end{equation}
The elementary symmetric functions are recovered recursively from the power sums:
\begin{equation}
 e_0=1,\qquad
 e_j=\frac1j\sum_{k=1}^j(-1)^{k-1}e_{j-k}s_k
 \quad(1\le j\le r).
 \label{eq:newton}
\end{equation}
For completeness, define $E(t)=\prod_i(1+p_it)=\sum_j e_jt^j$.
Logarithmic differentiation gives
\[
 E'(t)=E(t)\sum_{k\ge1}(-1)^{k-1}s_kt^{k-1}.
\]
Equating coefficients proves \eqref{eq:newton}.

\begin{proposition}[Finite exact identifiability]
\label{prop:exact}
The first $r$ even cumulants of $\mu_a$, together with the known background
cumulants, determine $a\in\sr$ uniquely.  The same is true of the raw moments
through order $2r$.  The reconstruction is: decode $s_1,\ldots,s_r$, compute
\eqref{eq:newton}, form \eqref{eq:poly}, list its roots with multiplicity, and take
their nonnegative square roots in increasing order.
\end{proposition}
\begin{proof}
\Cref{prop:decoder} supplies the power sums.  Newton's identities supply all
coefficients of $P_a$.  A monic polynomial determines its complex root multiset,
and in this model all those roots are real and nonnegative.  Nonnegative square
root and ordering therefore determine $a$.  Cumulants through order $2r$ are
polynomials in the raw moments through that order, proving the final assertion.
\end{proof}
No division by the observed characteristic function is used.  In particular,
zeros of the known up characteristic function cause no obstruction to this
exact finite-dimensional decoder.

\subsection{Total variation controls the required finite data}
Suppose temporarily that $a_r,b_r\le A$ and set $R=\max(1,L+rA)$.
All observed laws are supported on $[-R,R]$.  If
$\varepsilon=\TV(\mu_a,\mu_b)$, then their raw moments satisfy
\begin{equation}
 |m_n(\mu_a)-m_n(\mu_b)|\le2R^n\varepsilon.
 \label{eq:raw-bound}
\end{equation}

\begin{lemma}[An explicit cumulant bound]
\label{lem:cumulant-bound}
Let
$F_n=\sum_{q=1}^n q!\,\stirling{n}{q}$, where $\stirling{n}{q}$ is a Stirling number
of the second kind.  Then
\begin{equation}
 |\kap_n(\mu_a)-\kap_n(\mu_b)|
 \le2R^nF_n\varepsilon.
 \label{eq:cumulant-bound}
\end{equation}
Consequently
\begin{equation}
 \err(a,b)\le C_{r,R}\varepsilon,\qquad
 C_{r,R}=\max_{1\le k\le r}\frac{2R^{2k}F_{2k}}{|\gamma_{2k}|}.
 \label{eq:E-TV}
\end{equation}
The same bound holds through any other fixed finite cumulant order.
\end{lemma}
\begin{proof}
The moment--cumulant formula is
\[
 \kap_n(\mu)=\sum_{\pi\in\Pi_n}
 (-1)^{|\pi|-1}(|\pi|-1)!\prod_{D\in\pi}m_{|D|}(\mu),
\]
where $\Pi_n$ is the set of partitions of $\{1,\ldots,n\}$.
For a partition with $q$ blocks, telescope the difference of the two products.
Each of its $q$ terms is bounded by $2R^n\varepsilon$, by
\eqref{eq:raw-bound} and $|m_j|\le R^j$.  Multiplication by $(q-1)!$ and summation
over the $\stirling{n}{q}$ partitions gives \eqref{eq:cumulant-bound}.  The known
background cancels in differences in \eqref{eq:decoder}; division by the nonzero
$\gamma_{2k}$ gives \eqref{eq:E-TV}.
\end{proof}
These constants are explicit and conservative, not asserted to be numerically
optimal.  Their purpose is to make every finite-order continuity assertion
quantitative.

\subsection{Power-sum error controls polynomial-coefficient error}
For monic degree-$r$ polynomials, let $\norm{P-Q}_{\coef}$ be the largest absolute
difference between their nonleading coefficients.  Put $H_0=A^2$.
If $\delta=\err(a,b)$, then
\begin{equation}
 \norm{P_a-P_b}_{\coef}\le K_{r,A}\delta.
 \label{eq:coef-bound}
\end{equation}
One explicit choice is $K_{r,A}=\max_{1\le j\le r}K_j$, where
\begin{equation}
 K_0=0,\qquad
 K_j=\frac1j\sum_{k=1}^j
 \left(rH_0^kK_{j-k}+\binom r{j-k}H_0^{j-k}\right).
 \label{eq:K}
\end{equation}
Indeed, $|s_k|\le rH_0^k$ and $|e_j|\le\binom rjH_0^j$.  Subtracting the two
Newton recurrences, telescoping each product $e_{j-k}s_k$, and inducting on $j$
proves \eqref{eq:coef-bound} with exactly \eqref{eq:K}.

\subsection{The regular locus and its Jacobian}
The power-sum map in squared coordinates has Jacobian
\begin{equation}
 \frac{\partial s_k}{\partial p_i}=kp_i^{k-1},\qquad
 \det\left(\frac{\partial s_k}{\partial p_i}\right)_{k,i=1}^r
 =r!\prod_{i<j}(p_j-p_i).
 \label{eq:jacobian-p}
\end{equation}
The determinant follows by factoring $k$ from row $k$ and using the Vandermonde
determinant.  Passing to amplitude coordinates gives
\begin{equation}
 \det\left(\frac{\partial s_k}{\partial a_i}\right)_{k,i=1}^r
 =2^rr!\left(\prod_i a_i\right)\prod_{i<j}(a_j^2-a_i^2).
 \label{eq:jacobian-a}
\end{equation}
Thus the regular positive locus is precisely the locus of distinct positive
scales.  The determinant detects where local Lipschitz inversion can fail, but
by itself does not give the optimal exponent on the singular locus.  That is
the additional content of the main theorems.

\section{Root matching with multiplicities and separated cluster factors}
\label{sec:roots}

\subsection{A quantitative root-matching lemma}
A bound saying that each root is close to \emph{some} root of the other polynomial
is insufficient: it can ignore multiplicities.  We need a matching that counts
every root.

\begin{lemma}[Root matching on a bounded real interval]
\label{lem:root-match}
Let $P,Q$ be monic degree-$n$ polynomials with all roots in $[0,H_0]$.
There is a constant $C_{n,H_0}$ such that their ordered root lists $p,q$ satisfy
\begin{equation}
 \max_i|p_i-q_i|\le C_{n,H_0}\norm{P-Q}_{\coef}^{1/n}.
 \label{eq:root-match}
\end{equation}
For sufficiently small $\eta=\norm{P-Q}_{\coef}>0$, one may take
$C_{n,H_0}=4n(2H)^{1/n}$, with
$H=\sum_{\ell=0}^{n-1}(H_0+1)^\ell$.
\end{lemma}
\begin{proof}
For $|z|\le H_0+1$,
$|P(z)-Q(z)|\le H\eta$.  Set $\rho=(2H\eta)^{1/n}$ and suppose $\rho<1/2$.
Take the union of the disks of radius $\rho'$ centered at the roots of $P$, where
$\rho\le\rho'<2\rho$ is chosen to avoid tangencies of boundary circles.  All
boundary points satisfy $|z|\le H_0+1$ and have distance at least $\rho'$ from
every root of $P$.  Hence on that boundary,
\[
 |P(z)|=\prod_i|z-p_i|\ge(\rho')^n\ge2H\eta>|P(z)-Q(z)|.
\]
Rouch\'e's theorem, applied to the oriented boundary of each component, gives
the same number of zeros of $P$ and $Q$ in that component, with multiplicity.
Piecewise circular boundaries can equivalently be perturbed slightly before
applying the argument principle.  The total count is $n$, so every root of $Q$
is in one of these components.

A component contains at most $n$ disks and its disk-intersection graph is
connected.  Joining centers along this graph shows that its diameter is at most
$2n\rho'<4n\rho$.  Match equal numbers of roots within each component.  This gives
a full matching with displacement at most $4n\rho$.  Sorting the two real lists
cannot worsen the largest displacement, proving the displayed bound for small
$\eta$.  If $\eta=0$, the polynomials coincide.  For the remaining $\eta$, enlarge
the constant using the trivial bound $\max_i|p_i-q_i|\le H_0$ and the fixed
positive threshold for $\eta$.
\end{proof}

\subsection{Separated clusters depend Lipschitz-continuously on the coefficients}
\begin{lemma}[Lipschitz extraction of an isolated factor]
\label{lem:cluster}
Let $P^{\circ}$ be monic with distinct root clusters separated by disjoint
circles $\Gamma_0,\ldots,\Gamma_s$.  For nearby monic polynomials whose roots
remain inside the corresponding circles with the same cluster counts, the
monic factor belonging to any one circle is Lipschitz in the full coefficient
vector.  Its degree is the number of roots in that circle, counted with
multiplicity.
\end{lemma}
\begin{proof}
For a circle $\Gamma_j$ and $k\ge1$, the cluster power sum is
\begin{equation}
 t_{j,k}(P)=\frac1{2\pi i}\oint_{\Gamma_j}z^k\frac{P'(z)}{P(z)}\dd z.
 \label{eq:cluster-integral}
\end{equation}
The residue at a root of multiplicity $m$ is $m$ times the corresponding $k$th
power, so the formula includes all multiplicities.  Shrink the coefficient
neighborhood so that $|P|$ is uniformly bounded below on the fixed circles.
The numerators and derivatives are uniformly bounded there.  The identity
\[
 \frac{P'}P-\frac{Q'}Q=\frac{P'Q-Q'P}{PQ}
\]
then bounds the integrand difference by a constant times
$\norm{P-Q}_{\coef}$.  Integration gives a Lipschitz bound for the finitely many
cluster sums $t_{j,1},\ldots,t_{j,m_j}$.  Applying the finite Newton recurrences
to these sums gives the coefficients of the monic degree-$m_j$ factor and
preserves Lipschitz dependence on this bounded neighborhood.
\end{proof}
The use of a fixed contour is important.  The roots \emph{within} a cluster may
collide, but the cluster as a whole remains separated from the other clusters.
We are not assuming differentiable or Lipschitz individual roots at a collision.

\section{Upper bounds: proof of the pairwise and global estimates}
\label{sec:upper}

\begin{proof}[Proof of the upper bound in \cref{thm:pairwise}]
Choose small disjoint disks around $0$ when $m_0>0$ and around each positive
$c_j$.  Choose a relative parameter neighborhood $\mathcal U$ of $\base$ so that
exactly the prescribed coordinate block stays in each disk.  All parameters
in this neighborhood are bounded by some $A$.

Let $\delta=\err(a,b)$.  By \eqref{eq:coef-bound}, the coefficient vectors of
$P_a$ and $P_b$ differ by at most a constant times $\delta$.
\Cref{lem:cluster} gives the same type of estimate for the monic factor belonging
to cluster $j$.  Applying \cref{lem:root-match} to that degree-$m_j$ factor yields
\begin{equation}
 \max_{i\text{ in cluster }j}|a_i^2-b_i^2|
 \le C_j\delta^{1/m_j}.
 \label{eq:squared-cluster}
\end{equation}
Near a positive center $c_j$, both square roots are bounded away from zero.
The identity
\[
 |a_i-b_i|=\frac{|a_i^2-b_i^2|}{a_i+b_i}
\]
therefore preserves the exponent $1/m_j$.

For the zero cluster, the same argument in squared coordinates gives exponent
$1/m_0$.  For nonnegative $u,v$,
\begin{equation}
 |\sqrt u-\sqrt v|\le\sqrt{|u-v|},
 \label{eq:sqrt}
\end{equation}
so that cluster gives amplitude exponent $1/(2m_0)$.  Taking the maximum of
these bounds, and first treating $0\le\delta\le1$, gives exponent $1/M(\base)$.
An enlargement of the constant handles the remaining bounded range of $\delta$.
Finally, \cref{lem:cumulant-bound} gives
$\delta\le C_{r,R}\TV(\mu_a,\mu_b)$.
\end{proof}

\begin{proof}[Proof of the upper estimates in \cref{thm:global}]
On the full bounded parameter set, \eqref{eq:coef-bound} followed by
\cref{lem:root-match} in degree $r$ yields squared-scale exponent $1/r$.
Applying \eqref{eq:sqrt} gives amplitude exponent $1/(2r)$, and
\eqref{eq:E-TV} converts to total variation.  If all amplitudes are at least
$a_*>0$, the denominator $a_i+b_i\ge2a_*$ preserves exponent $1/r$ instead.

For a uniform squared-scale separation $\Delta>0$, sufficiently small
coefficient perturbations place each matched root $q_i$ within $\Delta/2$ of
$p_i$.  This follows uniformly from \cref{lem:root-match}.  Since $Q(q_i)=0$,
\[
 |P(q_i)|=|P(q_i)-Q(q_i)|\le H_1\norm{P-Q}_{\coef}
\]
for a fixed $H_1$ on the bounded interval.  Every factor
$|q_i-p_j|$, $j\ne i$, is at least $\Delta/2$, giving
\begin{equation}
 |q_i-p_i|\le H_1(2/\Delta)^{r-1}\norm{P-Q}_{\coef}.
 \label{eq:separated}
\end{equation}
Take square roots using the positive lower bound $a_*$, then use
\eqref{eq:coef-bound} and \eqref{eq:E-TV}.  Large coefficient perturbations are
covered by enlarging the constant, using the bounded parameter diameter and the
fixed threshold separating the small and large cases.  For $r=1$, the empty
product in \eqref{eq:separated} is one and no separation condition is needed.
\end{proof}

\begin{remark}[How much background regularity is used so far?]
The upper estimates use only that the known background law has bounded support.
They do not require a density or smoothness.  Smoothness becomes essential to
the sharpness constructions below, where high-order cancellation must persist
in the strong $L^1$ metric.  Thus the stated classification is a theorem for the
smooth-background model, while its upper half applies more widely.
\end{remark}

\section{Explicit moment-canceling clusters}
\label{sec:chebyshev}

The sharpness proof must stay inside the real nonnegative scale model.  Generic
perturbations of a polynomial with a repeated root often produce complex roots,
so they would not suffice.  Chebyshev level sets give suitable real configurations
at every multiplicity.

\begin{lemma}[Two real multisets with maximal initial cancellation]
\label{lem:cheb}
For $m\ge1$, let $x_1^+<\cdots<x_m^+$ and $x_1^-<\cdots<x_m^-$ be the roots of
$T_m(x)=1/2$ and $T_m(x)=-1/2$, respectively.  Both lists consist of simple roots
in $(-1,1)$ and have no root in common.  Their power sums satisfy
\begin{align}
 \sum_i(x_i^+)^k&=\sum_i(x_i^-)^k\quad(1\le k<m),
 \label{eq:cheb-cancel}\\
 D_m:=\sum_i(x_i^+)^m-\sum_i(x_i^-)^m&=m\,2^{1-m}>0.
 \label{eq:Dm}
\end{align}
The shifted lists $u_i^\pm=2+x_i^\pm\in(1,3)$ obey the same initial equalities
and the same $m$th-power difference $D_m$.
\end{lemma}
\begin{proof}
The identity $T_m(\cos\theta)=\cos(m\theta)$ shows that each horizontal level
in $(-1,1)$ meets the graph in exactly $m$ simple points of $(-1,1)$.  Simplicity
also follows because at such a point $\sin(m\theta)\ne0$ and
$T_m'(\cos\theta)=m\sin(m\theta)/\sin\theta\ne0$.

The monic polynomials of the two root lists are
\begin{equation}
 Q_\pm(z)=2^{1-m}\bigl(T_m(z)\mp1/2\bigr).
 \label{eq:cheb-poly}
\end{equation}
They have identical nonconstant coefficients.  The constant coefficient of
$Q_+$ minus that of $Q_-$ equals $-2^{1-m}$.  Newton's recurrence in coefficient form
says, for $1\le k\le m$,
\[
 s_k+c_1s_{k-1}+\cdots+c_{k-1}s_1+kc_k=0,
\]
where $c_j$ is the coefficient of $z^{m-j}$.  Thus the sums agree for $k<m$, and
at $k=m$ their difference is $-m(-2^{1-m})=D_m$.  The root sets are disjoint
because a number cannot simultaneously satisfy $T_m(x)=1/2$ and $T_m(x)=-1/2$.
Finally, expanding $(2+x)^k$ binomially proves both assertions for $u_i^\pm$.
\end{proof}

\begin{example}[A triple collision]
For $m=3$, the two monic polynomials are
\[
 Q_+(z)=z^3-\frac34z-\frac18,\qquad
 Q_-(z)=z^3-\frac34z+\frac18.
\]
Both root lists have sum zero and sum of squares $3/2$, but their sums of cubes
differ by $3/4$.  This produces two nearby three-factor configurations with
exactly the same second and fourth output cumulants, while the first differing
cumulant is sixth order.  A perturbation near a positive squared scale makes
that difference order $t^3$; a perturbation from zero in amplitude coordinates
makes it order $t^6$.
\end{example}

\section{Smoothness, cancellation, and leading density terms}
\label{sec:smoothness}

\subsection{Translation Taylor expansion in \texorpdfstring{$L^1$}{L1}}
\begin{lemma}[A uniform translation expansion]
\label{lem:translation}
Let $k\in C_c^\infty(\R)$ and let $Z$ be a bounded random variable.  For every
integer $n\ge1$,
\begin{equation}
 \E k(\,\cdot-tZ)
 =\sum_{\ell=0}^n\frac{(-t)^\ell\E Z^\ell}{\ell!}k^{(\ell)}
   +o_{L^1}(t^n)\quad(t\to0).
 \label{eq:translation}
\end{equation}
If two bounded variables have identical moments through order $n-1$, the
difference of their expansions starts with their $n$th-moment difference.
\end{lemma}
\begin{proof}
The integral Taylor formula for translation is
\[
 k(\,\cdot-h)-\sum_{\ell=0}^{n-1}\frac{(-h)^\ell}{\ell!}k^{(\ell)}
 =\frac{(-h)^n}{(n-1)!}\int_0^1(1-v)^{n-1}k^{(n)}(\,\cdot-vh)\dd v.
\]
Translations are continuous in $L^1$.  Replacing the translated $k^{(n)}$ in the
integral by $k^{(n)}$ therefore changes the expression by $o_{L^1}(|h|^n)$,
uniformly for $|h|$ in a shrinking interval.  Substitute $h=tZ$.  Boundedness of
$Z$ makes the remainder uniform over its possible values; taking expectations
proves \eqref{eq:translation}.
\end{proof}

Convolving $g$ with any fixed collection of finite uniform factors gives another
smooth compactly supported density $k$, with
$\norm{k^{(n)}}_1\le\norm{g^{(n)}}_1$.  Moreover,
\begin{equation}
 \norm{k^{(n)}}_1>0\qquad(n\ge1).
 \label{eq:nonzero-derivative}
\end{equation}
Otherwise $k^{(n)}=0$ everywhere by continuity, so $k$ is a polynomial of degree
less than $n$ on $\R$.  A compactly supported polynomial vanishes identically,
contradicting $\int k=1$.

\subsection{Smooth symmetric dependence on positive squared scales}
\begin{lemma}[Positive-cluster Taylor cancellation]
\label{lem:positive-Taylor}
For a fixed smooth compactly supported density $k$ and a fixed $c>0$, define
\[
 \mathcal F(p_1,\ldots,p_m)(x)
 =\E k\left(x-\sum_{i=1}^m\sqrt{p_i}U_i\right)
 \quad(p_i>0).
\]
This is a symmetric $C^\infty$ map into $L^1(\R)$.  If two direction vectors
$v,w\in\R^m$ have the same power sums through order $m-1$, then
\begin{equation}
 \mathcal F(c\mathbf1+tv)-\mathcal F(c\mathbf1+tw)
 =t^m H+o_{L^1}(t^m)
 \label{eq:positive-expansion}
\end{equation}
for some $H\in L^1$.
\end{lemma}
\begin{proof}
On a compact subset of the positive orthant, every derivative of $\sqrt p$ is
bounded.  Differentiating the displayed expectation produces finite sums of
bounded functions of the $U_i$ times translated derivatives $k^{(j)}$.  Their
$L^1$ norms have uniform integrable bounds.  Translation continuity and
differentiation under the expectation give $C^\infty$ dependence into $L^1$.
Permutation symmetry follows from the identical independent laws of the $U_i$.

Expand the map at $c\mathbf1$ through total degree $m$.  Its homogeneous term of
degree $\ell$ is a symmetric polynomial in the direction coordinates, with
coefficients in $L^1$.  Every symmetric polynomial of degree $\ell$ is a
polynomial in the elementary symmetric functions $e_1,\ldots,e_\ell$ with their
usual weights, hence in the power sums $s_1,\ldots,s_\ell$ by Newton's identities.
This algebraic statement also holds for vector-valued coefficients: their span
in this finite Taylor polynomial is finite dimensional, and one applies the
scalar identity in a basis of that span.  One elementary proof of the scalar
statement successively subtracts the product
$e_1^{\lambda_1-\lambda_2}e_2^{\lambda_2-\lambda_3}\cdots e_m^{\lambda_m}$
corresponding to the leading ordered monomial $z_1^{\lambda_1}\cdots z_m^{\lambda_m}$;
the leading monomial decreases until none remains.  The weighted degree is
preserved throughout.

Consequently all Taylor terms of degrees below $m$ agree for $v$ and $w$.
The difference of their degree-$m$ terms is $H$, and the Banach-space Taylor
remainder is $o_{L^1}(t^m)$, proving \eqref{eq:positive-expansion}.
\end{proof}

\section{Sharpness at positive collisions}
\label{sec:positive-sharp}

Fix a positive squared-scale cluster of size $m$ at $c>0$, and hold all other
coordinates equal to their values at $\base$.  Let $k$ be the density of the
known background convolved with these fixed factors.  For the Chebyshev lists
from \cref{lem:cheb}, replace the cluster by
\begin{equation}
 p_i^\pm(t)=c+t x_i^\pm,\qquad
 a_i^\pm(t)=\sqrt{c+t x_i^\pm},\qquad t>0.
 \label{eq:positive-paths}
\end{equation}
For small $t$, all scales remain positive and in their original isolated cluster.
The full configurations converge to $\base$ and satisfy
\begin{equation}
 d(a^+(t),a^-(t))\sim
 \frac{t}{2\sqrt c}\max_i|x_i^+-x_i^-|.
 \label{eq:positive-distance}
\end{equation}
The coefficient on the right is positive because the two ordered root lists
differ.

By binomial expansion and \cref{lem:cheb}, the cluster power sums agree through
order $m-1$, while their order-$m$ difference is exactly $D_mt^m$.
The fixed coordinates do not change those differences.  Therefore
\begin{equation}
 \kap_{2m}(\mu_{a^+(t)})-\kap_{2m}(\mu_{a^-(t)})
 =\gamma_{2m}D_mt^m\ne0,
 \label{eq:positive-kappa}
\end{equation}
and all lower-order cumulants of the output laws agree.  Odd cumulants, even
when the background is not symmetric, agree because the uniform factors have
zero odd cumulants and the background is the same.

\begin{proposition}[Positive-collision sharpness in total variation]
\label{prop:positive-sharp}
For the paths \eqref{eq:positive-paths}, there is a nonzero compactly supported
$H_{c,m,k}\in L^1(\R)$ such that
\begin{align}
 f_{a^+(t)}-f_{a^-(t)}
 &=t^mH_{c,m,k}+o_{L^1}(t^m),
 \label{eq:positive-density}\\
 \TV(\mu_{a^+(t)},\mu_{a^-(t)})
 &\sim\frac12\norm{H_{c,m,k}}_1t^m.
 \label{eq:positive-TV}
\end{align}
In addition, $\err(a^+(t),a^-(t))\asymp t^m$.
\end{proposition}
\begin{proof}
\Cref{lem:positive-Taylor} gives the first expansion.  All densities along these
paths have support in one fixed compact interval; the Taylor coefficients,
being derivatives of such densities, have support there as well.
Since all lower cumulants agree, the moment--cumulant polynomial implies that
the raw moments of order $2m$ differ by the cumulant difference
\eqref{eq:positive-kappa}.  Dividing the $L^1$ expansion by $t^m$ and integrating
against $x^{2m}$ on that fixed compact interval gives
\[
 \int x^{2m}H_{c,m,k}(x)\dd x=\gamma_{2m}D_m\ne0.
\]
Thus $H_{c,m,k}$ is nonzero.  Taking $L^1$ norms proves
\eqref{eq:positive-TV}.  Finally, the $m$th power-sum gap is $D_mt^m$, while
for every fixed $k\ge m$ the binomial expansion makes its gap $O(t^m)$.
As $m\le r$, the maximum over $1\le k\le r$ is bounded above and below by
positive constants times $t^m$.
\end{proof}

By \eqref{eq:positive-distance} and \eqref{eq:positive-TV}, a putative pairwise
bound with exponent $\alpha>1/m$ would require
$t\le C't^{m\alpha}$ along these paths, which is impossible as $t\downarrow0$.
This proves the positive-cluster sharpness half of \cref{thm:pairwise}.
The argument includes $m=1$ and therefore also proves that a regular exponent
larger than one cannot hold.

\section{Vanishing factors: an exact leading law and doubled exponents}
\label{sec:zero-sharp}

The zero cluster has an extra degeneracy: the law of $aU$ is even as a function
of the amplitude $a$, so its first detectable change occurs in $a^2$, not $a$.
The next theorem quantifies how this interacts with an $m$-fold collision.

\begin{theorem}[Exact vanishing-cluster leading term]
\label{thm:zero-leading}
Suppose $\base$ has a zero cluster of size $m\ge1$.  Hold all other scales fixed,
and let $k$ be the smooth compactly supported density of the background plus
those fixed factors.  Set $u_i^\pm=2+x_i^\pm$ using \cref{lem:cheb}, and turn the
zero cluster into
\begin{equation}
 a_i^\pm(t)=t\sqrt{u_i^\pm},\qquad t>0.
 \label{eq:zero-paths}
\end{equation}
Then
\begin{equation}
 f_{a^+(t)}-f_{a^-(t)}
 =\frac{\gamma_{2m}D_m}{(2m)!}\,t^{2m}k^{(2m)}
   +o_{L^1}(t^{2m}),
 \label{eq:zero-leading}
\end{equation}
and in particular
\begin{equation}
 \TV(\mu_{a^+(t)},\mu_{a^-(t)})
 \sim\frac{|\gamma_{2m}D_m|}{2(2m)!}
       \norm{k^{(2m)}}_1\,t^{2m}.
 \label{eq:zero-TV}
\end{equation}
The constant is strictly positive.  Meanwhile,
\begin{equation}
 d(a^+(t),a^-(t))
 =t\max_i|\sqrt{u_i^+}-\sqrt{u_i^-}|>0,
 \label{eq:zero-distance}
\end{equation}
and $\err(a^+(t),a^-(t))\asymp t^{2m}$.
\end{theorem}
\begin{proof}
Define bounded symmetric random variables
$Z_\pm=\sum_i\sqrt{u_i^\pm}U_i$.  Their even cumulants are
$\gamma_{2j}\sum_i(u_i^\pm)^j$.  By \cref{lem:cheb}, these agree for $j<m$,
and their order-$2m$ difference is $\gamma_{2m}D_m$.  Odd cumulants vanish.
The moment--cumulant identities therefore give equal raw moments through order
$2m-1$ and
\[
 \E Z_+^{2m}-\E Z_-^{2m}=\gamma_{2m}D_m.
\]
The two output densities are $\E k(\cdot-tZ_+)$ and $\E k(\cdot-tZ_-)$.
Apply \cref{lem:translation} with $n=2m$ and subtract.  The sign $(-1)^{2m}$ is
positive, giving exactly \eqref{eq:zero-leading}.

Both $\gamma_{2m}$ and $D_m$ are nonzero, and
$\norm{k^{(2m)}}_1>0$ by \eqref{eq:nonzero-derivative}.  Norm convergence gives
\eqref{eq:zero-TV}.  The amplitude formula is immediate.  In squared coordinates
$p_i^\pm=t^2u_i^\pm$, the sums agree through order $m-1$ and differ at order
$m$ by $D_mt^{2m}$.  Higher fixed orders have differences $O(t^{2m+2})$ or
smaller, proving the final assertion.
\end{proof}

For $m=3$, the signed coefficient in \eqref{eq:zero-leading} is
\[
 \frac{(16/63)(3/4)}{6!}=\frac1{3780}.
\]
Thus a vanishing triple has density difference
$t^6k^{(6)}/3780+o_{L^1}(t^6)$ and total variation asymptotic
$t^6\norm{k^{(6)}}_1/7560$.  The amplitude separation is proportional to $t$.
This is an explicit sixth-order obstruction, not merely a qualitative
non-Lipschitz example.

\begin{proof}[Completion of \cref{thm:pairwise,thm:global}]
Choose a cluster attaining $M(\base)$.  If it is positive, use
\cref{prop:positive-sharp}; if it is zero, use \cref{thm:zero-leading}.  In either
case the two configurations tend to $\base$, their amplitude separation is
comparable to $t$, and both their finite-data error and their total variation
are comparable to $t^{M(\base)}$.  Every neighborhood therefore contains pairs
excluding every exponent greater than $1/M(\base)$.  Together with
\cref{sec:upper}, this proves \cref{thm:pairwise} in full.

For the full box $0\le a_i\le A$, use a zero cluster of size $r$ at the origin.
For the positive box $a_*\le a_i\le A$, choose all coordinates equal to any
value strictly between $a_*$ and $A$, and use a positive cluster of size $r$.
These constructions remain inside their boxes for small $t$ and prove the
claimed optimal global exponents in \cref{thm:global}.
\end{proof}

\section{Anchored recovery and why its exponent is different}
\label{sec:anchored}

\subsection{A nonnegative polynomial certificate}
The poor pairwise examples compare two moving, highly correlated clusters.
When one configuration is held exactly at the singular base, a nonnegative
polynomial blocks this cancellation.

\begin{proof}[Upper bound in \cref{thm:anchored}]
The simple positive case follows from \cref{thm:pairwise}.  For the other cases,
let $c_1,\ldots,c_s$ be the distinct positive squared centers of the base.
Define
\begin{equation}
 Q(x)=
 \begin{cases}
 \displaystyle\prod_{j=1}^s(x-c_j)^2,&m_0=0,\\[0.4em]
 \displaystyle x\prod_{j=1}^s(x-c_j)^2,&m_0>0.
 \end{cases}
 \label{eq:anchor-certificate}
\end{equation}
The empty product is one, so if all coordinates vanish then $Q(x)=x$.
This polynomial is nonnegative for $x\ge0$ and vanishes at every base squared
coordinate.  If $Q(x)=\sum_{\ell=0}^d q_\ell x^\ell$, then
\begin{equation}
 0\le\sum_i Q(a_i^2)
 =\sum_{\ell=1}^d q_\ell\bigl(s_\ell(a)-s_\ell(\base)\bigr)
 \le C\TV(\mu_a,\mu_{\base}).
 \label{eq:anchor-certificate-bound}
\end{equation}
The constant terms cancel because both lists have length $r$.  The last
inequality uses the decoder and \cref{lem:cumulant-bound} through the fixed
finite order $2d$; it does not assume $d\le r$.

Near a positive center $c_j$, the ratio $Q(x)/(x-c_j)^2$ is continuous and
strictly positive.  Hence $Q(x)\ge b_j(x-c_j)^2$ there for some $b_j>0$.
When there is a zero cluster, $Q(x)/x=\prod_j(x-c_j)^2$ is strictly positive
near zero, so $Q(x)\ge b_0x$ for small $x\ge0$.  Shrink the parameter
neighborhood so its coordinate blocks stay near their prescribed centers.
Every summand in \eqref{eq:anchor-certificate-bound} is nonnegative.
It follows that
\[
 |a_i^2-c_j|\le C\sqrt\varepsilon\quad\text{in a positive block},
 \qquad a_i^2\le C\varepsilon\quad\text{in the zero block},
\]
where $\varepsilon=\TV(\mu_a,\mu_{\base})$.  The positive-block square-root map
is Lipschitz, while in the zero block taking square roots gives
$a_i\le C\sqrt\varepsilon$.  Combining the blocks proves
\eqref{eq:anchored}.
\end{proof}

\subsection{Optimality of the square-root anchored exponent}
\begin{proof}[Sharpness and completion of \cref{thm:anchored}]
If a zero coordinate is present, change just that amplitude from zero to $t$,
leaving all other coordinates fixed and sorting afterward.  Let $k$ be the density
of the background and those fixed factors.  Since $\E U=0$ and $\E U^2=1/3$,
\cref{lem:translation} gives
\begin{equation}
 f_t-k=\frac{t^2}{6}k''+o_{L^1}(t^2),\qquad
 \TV(\mu_t,\mu_{\base})\sim\frac{\norm{k''}_1}{12}t^2.
 \label{eq:anchored-zero}
\end{equation}
The parameter distance is $t$.  The positive leading constant excludes every
anchored exponent greater than $1/2$.

If there is no zero but a repeated positive squared scale $c$, replace two
copies of $c$ by $c+t$ and $c-t$.  The first cluster power sum stays equal to
$2c$, while the second increases by $2t^2$.  Positive-coordinate smoothness
shows that the density has an $L^1$ Taylor expansion.  Its linear term cancels
by symmetry in the two chosen coordinates, so the difference from the base is
$t^2H+o_{L^1}(t^2)$.  The output fourth cumulant changes by
$2\gamma_4t^2\ne0$, whereas all lower cumulants agree.  Integrating the expansion
against $x^4$ on a common compact support gives $\int x^4H=2\gamma_4\ne0$.
Thus total variation is comparable to $t^2$, while amplitude distance is
comparable to $t$.  Again an exponent greater than $1/2$ is impossible.

Finally, at a simple positive base, change one squared scale from $c$ to $c+t$.
Smooth dependence gives an $O(t)$ upper bound in total variation, and the
variance change $\gamma_2t$ gives a matching lower bound by
\cref{lem:cumulant-bound}.  Parameter distance is comparable to $t$, excluding
anchored exponents greater than one.  The previously proved upper bounds
complete all cases.
\end{proof}

\begin{remark}[There is no contradiction between the two classifications]
At an $m$-fold positive collision, each Chebyshev configuration is generally
at law distance of order $t^2$ from the singular base, while their distance
from each other can be of order $t^m$.  At a zero cluster, each configuration
has a variance change of order $t^2$ from the base, even though the two output
laws can differ only at order $t^{2m}$.  Anchoring removes the possibility that
their common low-order changes cancel against one another.
\end{remark}

\section{Deterministic recovery, feasible fitting, and sampling}
\label{sec:recovery}

\subsection{A forward continuity bound}
Using the same uniforms to couple two models and the translation bound
$\norm{g(\cdot-u)-g(\cdot-v)}_1\le\norm{g'}_1|u-v|$, we obtain
\begin{equation}
 \TV(\mu_a,\mu_b)
 \le\frac{\norm{g'}_1}{4}\sum_{i=1}^r|a_i-b_i|
 \le\frac{r\norm{g'}_1}{4}d(a,b).
 \label{eq:forward}
\end{equation}
Here $\E|U_i|=1/2$ supplies the factor $1/4$ after the total-variation
normalization.  This bound ensures continuity of the forward map on compact
parameter sets.  It need not be sharp at singular points.

\subsection{Optimal deterministic ambiguity radii}
Take a compact relative neighborhood $K$ of $\base$ small enough that
\eqref{eq:main-bound} holds.  Suppose an observed probability law $\nu$ obeys
$\TV(\nu,\mu_a)\le\varepsilon$ for an unknown $a\in K$.  Minimize
$\TV(\nu,\mu_b)$ over $b\in K$; continuity and compactness ensure a minimizer
$\widehat a$.  It satisfies
\[
 \TV(\mu_{\widehat a},\mu_a)\le2\varepsilon,
 \qquad d(\widehat a,a)\le C(2\varepsilon)^{1/M(\base)}.
\]
An exact minimizer is used to state the information-theoretic bound, not to
claim an efficient numerical algorithm.

Conversely, take either sharpness pair $\mu_+(t),\mu_-(t)$ and let
$\nu_t=(\mu_+(t)+\mu_-(t))/2$.  Each model is at total-variation distance
$\TV(\mu_+(t),\mu_-(t))/2$ from $\nu_t$.  Every single output of an estimator
on this common observation is at parameter distance at least
$d(a^+(t),a^-(t))/2$ from one of the two models.  With
$t=c\varepsilon^{1/M}$ and a sufficiently small fixed $c$, both models lie
in the allowed $\varepsilon$-ball for all sufficiently small $\varepsilon$.
Therefore the worst-case deterministic recovery error is bounded below by a
positive constant times $\varepsilon^{1/M}$.  Together with the upper bound,
this proves the optimal deterministic rate near the base.

\subsection{Exact finite-data reconstruction is not automatically noise-safe}
For exact data, the decoder in \cref{prop:exact} is sufficient.  Noisy data
require an additional feasibility step.  A polynomial formed by blindly
applying Newton's identities to perturbed sums may have negative or nonreal
roots, even when the perturbation is tiny.  Taking real parts of its roots is
not justified by the theorems in this article.

A mathematically safe alternative is the following finite-dimensional fit.
Given approximate sums $\widetilde s_1,\ldots,\widetilde s_r$ and a compact
parameter set $K$, minimize
\begin{equation}
 \rho(b)=\max_{1\le k\le r}|s_k(b)-\widetilde s_k|
 \quad\text{over }b\in K.
 \label{eq:feasible-fit}
\end{equation}
A minimizer exists.  If the true parameter $a\in K$ has
$\max_k|s_k(a)-\widetilde s_k|\le\eta$, then every minimizer $\widehat a$ has
\begin{equation}
 \err(\widehat a,a)\le2\eta,
 \qquad d(\widehat a,a)\le C(2\eta)^{1/M}
 \label{eq:moment-fit-bound}
\end{equation}
when $K$ is a local neighborhood as above.  On the global box, the corresponding
exponent is $1/(2r)$; on a positive box it is $1/r$.  An approximate minimizer
with objective at most $\inf\rho+\tau$ gives $2\eta+\tau$ in place of $2\eta$.

For certified computation one may work in squared coordinates
$0\le p_1\le\cdots\le p_r\le A^2$ and use interval subdivision or polynomial
inequality methods to enclose the feasible set.  This is a valid route to a
certificate, but a fast implementation and tight complexity analysis are not
provided here.  The included executable code validates the exact algebraic
stage, not a complete certified nonlinear optimizer.

\subsection{A sampling upper bound, with its limits made explicit}
The total variation distance between a continuous law and its finitely supported
empirical distribution is one.  Consequently the deterministic TV estimate
must \emph{not} be applied directly to empirical TV distance.  Bounded moments
provide a different route.

\begin{proposition}[Elementary sampling upper bound]
\label{prop:sampling}
Assume independent samples $Y_1,\ldots,Y_N$ from $\mu_a$, all supported on
$[-R,R]$, and a known compact parameter neighborhood $K$ as in
\eqref{eq:main-bound}.  There is a constant $C$ such that a feasible moment fit
satisfies, with probability at least $1-\beta$,
\begin{equation}
 d(\widehat a,a)\le
 C\left(\frac{\log(4r/\beta)}{N}\right)^{1/(2M(\base))},
 \label{eq:sampling}
\end{equation}
for sufficiently large $N$.  Without a known local neighborhood, the global
bounded-box version gives exponent $1/(4r)$ in $N$.
\end{proposition}
\begin{proof}
For each $1\le n\le2r$, the sample average of $Y_j^n$ has expectation
$m_n(\mu_a)$.  Normalize by $R^n$, so the summands lie in $[-1,1]$.
A bounded-variable exponential estimate gives
\[
 \Pr\left(|\widetilde m_n-m_n|>R^n u\right)
 \le2e^{-Nu^2/2}.
\]
For clarity, this follows by bounding the centered moment generating function
of a variable in an interval of length two by $e^{\lambda^2/2}$, applying
Markov's inequality to the sum, and optimizing at $\lambda=u$; apply the same
argument to the negative sum.  To justify the moment-generating-function bound,
let $h(\lambda)=\log\E e^{\lambda(W-\E W)}$.  Then $h(0)=h'(0)=0$ and
$h''(\lambda)$ is the variance under an exponential tilt.  A distribution
supported on an interval of length two has variance at most one, since its
mean minimizes mean squared distance and its distance from the interval
midpoint is at most one.  Thus $h''\le1$; integrating twice yields
$h(\lambda)\le\lambda^2/2$.  A union bound over $2r$ moments and the choice
$u=\sqrt{2\log(4r/\beta)/N}$ give simultaneous moment bounds.

The finite cumulant polynomials are Lipschitz on the bounded set containing
both the true and sample moment vectors.  Dividing by the fixed nonzero
$\gamma_{2k}$ and subtracting the known background cumulants gives sum error
$\eta\le C_0u$.  Apply \eqref{eq:moment-fit-bound}.  The global statement follows
from the global exponent $1/(2r)$ for finite-data error.
\end{proof}
We do not claim that \eqref{eq:sampling} is statistically minimax.  A sharp
statistical lower bound needs control of testing distances such as Hellinger
distance or relative entropy, not merely the TV expansions between smooth
densities established here.  Endpoint flatness can matter to that stronger
problem.

\subsection{Rank is not continuously detectable at zero}
Within capacity $r$, a zero factor can be replaced by an arbitrarily small
positive one, changing the law by order $t^2$ while changing the number of
positive factors by one.  Thus the exact positive-factor count is not a
continuous function of the observed law at a configuration with a zero slot.
A positive lower scale bound or a thresholded notion of effective rank is
necessary for robust rank reporting.  This does not contradict exact recovery
of the padded multiset from exact moments.

\section{Verification record and reproducible computations}
\label{sec:verification}

\subsection{Exact checks}
The accompanying \texttt{verification.py} was executed with Python 3.13.5,
SymPy 1.14.0, and mpmath 1.3.0.  The exact checks use rational polynomial
arithmetic.  For every $m=1,\ldots,12$, they verify the monic Chebyshev polynomial
pair, equality of power sums through $m-1$, the formula $D_m=m2^{1-m}$,
the corresponding shifted-list equalities, simplicity of each polynomial,
and disjointness of the two root sets.  These are finite tests of identities
already proved for all $m$ in \cref{lem:cheb}.

For capacities $r=1,\ldots,8$, exact examples containing zero, repeated, and
distinct rational squared scales were converted to output cumulants over the
up background.  The code then converts cumulants to raw moments, recovers
cumulants from those moments, decodes the power sums, and reconstructs the
monic polynomial.  Each reconstructed polynomial equals the original exactly.
For example, squared scales $(0,3/5,1/5,1/5)$ give
\begin{equation}
 P(z)=z^4-z^3+\frac7{25}z^2-\frac3{125}z.
 \label{eq:checked-poly}
\end{equation}
Sorting the scale list does not change this polynomial.

\subsection{High-precision Fourier diagnostics}
A separate numerical diagnostic evaluates characteristic-function differences
at the single frequency $\xi=0.9$.  It uses 100 decimal digits of precision,
Chebyshev clusters of orders $3$ and $5$, and $t=2^{-\ell}$ for
$\ell=3,\ldots,12$.  The fixed up multiplier is approximated by its first 220
sinc factors.  At this frequency the omitted logarithmic tail is less than
$10^{-132}$ in absolute value: for $|x|\le1/2$ one can bound
$|\log\sinc x|\le x^2$, and the resulting tail is a geometric sum.
This multiplier is common to the two configurations, so it does not control
the observed cancellation order.

If $G(t)$ is the absolute Fourier difference, the diagnostic reports
$\log_2(G(2t)/G(t))$.  Its last computed values are:
\begin{center}
\begin{tabular}{@{}lcrr@{}}
\toprule
Regime & Cluster size & Predicted power & Last halving slope\\
\midrule
Positive collision & 3 & 3 & 2.999999999868\\
Vanishing cluster & 3 & 6 & 5.999999916416\\
Positive collision & 5 & 5 & 4.999999999813\\
Vanishing cluster & 5 & 10 & 9.999999863296\\
\bottomrule
\end{tabular}
\end{center}
These are pointwise Fourier diagnostics, \emph{not} numerical evaluations of
total variation.  They do not prove the analytical theorems, their $L^1$
remainders, or the optimality statements.  Those conclusions follow from the
proofs above.  The explicit distinction prevents a cancellation plot at one
frequency from being mistaken for a strong-metric convergence certificate.

\subsection{How to reproduce and what was not checked}
From the unpacked directory, run
\begin{verbatim}
python -m pip install -r requirements.txt
python verification.py --output-dir .
pdflatex -interaction=nonstopmode -halt-on-error article.tex
pdflatex -interaction=nonstopmode -halt-on-error article.tex
pdflatex -interaction=nonstopmode -halt-on-error article.tex
\end{verbatim}
The script writes \texttt{verification\_results.json} and two CSV tables under
\texttt{data/}.  The assertions test exact identities and the stated numerical
slopes.  No infinite family is certified merely by passing a finite range,
and the Python program is not part of a theorem-prover trusted boundary.
No Lean or Rocq build was performed for this article.  In particular, the
repository's existing formal identifiability result must not be interpreted
as already formalizing the different stability theorems proved here.

\section{A concrete formalization roadmap}
\label{sec:formalization}

A useful implementation should keep the algebraic and analytic boundaries
separate.  The algebraic core is smaller than the full probability-theoretic
classification and admits independent testing.

\paragraph{Finite symmetric algebra.}
Formalize sorted nonnegative lists, squared-coordinate multisets, power sums,
Newton reconstruction, and the coefficient recurrence.  The initial endpoint
is equality of multisets from the first $r$ power sums over a characteristic-zero
field.  This should not depend on Fabius-specific analysis.

\paragraph{Quantitative root matching.}
Formalize an actual multiplicity-preserving bijection of roots, not a
Hausdorff-set approximation.  A complex-analytic route follows
\cref{lem:root-match}; a real-root-specific alternative might reduce dependencies.
The isolated-factor lemma can then be proved by logarithmic contour integrals
or an algebraic implicit-function argument at coprime factors.  Neither route
is represented here as already implemented.

\paragraph{Probability and finite cumulants.}
Connect the actual law of $B+\sum a_iU_i$ to finite raw moments and cumulants.
Prove the uniform coefficient formula, nonvanishing, the known-background
subtraction, and the explicit TV-to-moment bounds.  The known up construction
then becomes a specialization rather than part of every argument.

\paragraph{Smooth \texorpdfstring{$L^1$}{L1} parameter dependence.}
The hardest analytic bridge is a reusable translation Taylor theorem in $L^1$,
followed by parameter differentiation of finite uniform convolutions.  The
Banach-valued symmetric Taylor argument must be connected to actual densities,
not merely to formal characteristic-function series.

\paragraph{Two independent sharpness paths.}
Encode the Chebyshev level polynomials and their Newton sums.  For positive
collisions, identify a nonzero leading density coefficient using a compactly
supported moment functional.  For zero clusters, prove the exact coefficient
in \cref{thm:zero-leading}.  These produce genuine counterexamples to stronger
exponents.  A proof of an upper bound alone is not a formalization of optimality.

\paragraph{Anchored positivity certificates.}
The polynomial $Q$ in \eqref{eq:anchor-certificate} gives a separate, comparatively
elementary route to fixed-base stability.  It is a good intermediate target:
it captures the conceptual difference between anchored and pairwise inversion
without requiring the entire Chebyshev family.

The desired final ledger should distinguish exact algebraic proofs, analytic
bridges, imported background results, executable numerical tests, and remaining
obligations.  This article supplies theorem statements and conventional proofs,
not untested code presented as machine-checked mathematics.

\section{Further research questions and proposed projects}
\label{sec:questions}

The following are proposed directions beyond the theorems proved here.  Their
status is a research agenda, not a verified claim that each is a previously
published open problem.

\begin{problem}[The optimal statistical experiment near collisions]
For the same smooth-background model, determine the local minimax rate from
$N$ independent samples at each collision pattern.  Does the rate
$N^{-1/(2M)}$ from the moment upper bound hold sharply under natural conditions
on the background?  The concrete first target is to compute the Hellinger
asymptotic of the Chebyshev pairs, including the contribution near the support
boundary.  An $L^1$ expansion alone does not justify squaring its leading term
and dividing by a possibly flat density.
\end{problem}

\begin{problem}[The exact smoothing threshold]
Replace $C_c^\infty$ by finite Sobolev regularity, or remove the background
altogether.  Determine which collision patterns still have TV exponent $1/M$.
The upper bounds already hold for any compactly supported known background.
The missing issue is whether insufficient regularity exposes low-order moving
knots before the moment-canceling mechanism reaches order $m$ or $2m$.
\end{problem}

\begin{problem}[Unbounded capacity and tail-controlled infinite products]
Let the unknown scale multiset be infinite with a prescribed summability or
geometric-tail bound.  Identify a metric on scale multisets and the sharp
modulus of recovery.  The finite-capacity exponent $1/(2r)$ degenerates with $r$,
so no rank-independent positive power law can be inferred from this article.
A useful first problem is to combine a certified tail truncation bound with the
finite theorem and optimize the truncation depth as a function of noise.
\end{problem}

\begin{problem}[Unknown background parameters]
Allow the background to belong to a known finite-dimensional family rather
than being fixed.  Determine when its cumulant directions can be separated
from the unknown uniform-factor directions.  Merely subtracting an estimated
background is not enough if the two parameter families share tangent or
higher-order cumulant directions.  A concrete target is a known smooth kernel
with an additional unknown scale and a finite collection of uniform factors.
\end{problem}

\begin{problem}[Certified feasible fitting with sharp constants]
Design an efficient certified algorithm for \eqref{eq:feasible-fit} or for
its interval-data feasibility version.  Quantify bit complexity in capacity,
separation, requested amplitude precision, and distance to the nonnegative
boundary.  The theorem supplies optimal exponents, but the explicit general
constants in \cref{lem:cumulant-bound} are deliberately coarse.
\end{problem}

\begin{problem}[Arithmetic scale restrictions]
Study how the inverse problem changes when scales lie on a known dyadic or
other discrete geometric grid.  The repository's analytic-order decoder
suggests a different source of rigidity from unrestricted real-scale moment
inversion.  Determine noise thresholds for exact grid-scale recovery and compare
zero-divisor measurements with finite cumulant measurements on the same model.
\end{problem}

\begin{problem}[Multivariate segment convolutions]
Replace scalar amplitudes by unknown vectors and consider sums of uniform
segments in $\R^d$ over a known smooth background.  Even cumulants become
symmetric tensors.  Classify identifiability and stability modulo permutation
and segment reversal, distinguishing collinear collisions from genuinely
multidimensional degeneracies.  The scalar Chebyshev obstruction embeds in every
dimension, but need not exhaust the singularities.
\end{problem}

\begin{problem}[First complete formal classification]
Formalize both \cref{thm:pairwise,thm:anchored}, including explicit sharpness
paths, in a proof assistant.  An attractive staged project is first exact
finite recovery, then anchored positivity certificates, then global root
matching, and finally the smooth $L^1$ sharpness arguments.
\end{problem}

\section{Conclusion}
The finite unknown-factor problem has a precise answer once its metric,
capacity, background, and symmetries are fixed.  Cumulants remove the known
background and turn the unknown factors into a real-root polynomial.  Positive
root collisions determine one family of exponents; the square-root map at the
nonnegative boundary doubles the zero-cluster multiplicity.  Smoothness ensures
that explicit moment cancellation survives in total variation, making the
result sharp for actual probability laws.

The same analysis shows why the quantifiers matter.  Uniform pairwise recovery
near a singular configuration has exponent $1/M$, whereas recovery relative
to that exact base has exponent $1/2$, except at the regular positive locus where
it is Lipschitz.  Confusing these two questions would obscure the strongest
instability and simultaneously underestimate anchored stability.

The result is a finite-rank stability theorem with complete conventional proofs
and reproducible algebraic checks.  Infinite capacity, optimal statistical
rates, weaker smoothing, and a full formal implementation remain separate
research tasks rather than consequences asserted without proof.

\begin{thebibliography}{9}
\raggedright
\bibitem{Arias}
J.~Arias de Reyna,
\emph{An infinitely differentiable function with compact support: Definition
and properties}, arXiv:1702.05442 (2017), English translation of the 1982 paper
in \emph{Rev. Real Acad. Ciencias Madrid} \textbf{76}, 21--38.
\url{https://arxiv.org/abs/1702.05442}.

\bibitem{Prony}
D.~Batenkov and Y.~Yomdin,
\emph{Geometry and Singularities of the Prony mapping},
Journal of Singularities \textbf{10} (2014), 1--25.
DOI: \href{https://doi.org/10.5427/jsing.2014.10a}{10.5427/jsing.2014.10a}.
\url{https://journalofsingularities.org/volume10/article1.html}.

\bibitem{Schaetzle}
R.~Sch\"atzle,
\emph{On the perturbation of the zeros of complex polynomials},
IMA Journal of Numerical Analysis \textbf{20}(2) (2000), 185--202.
DOI: \href{https://doi.org/10.1093/imanum/20.2.185}{10.1093/imanum/20.2.185}.

\bibitem{RepoIdent}
V.~Reshetnikov, \emph{ProveIt},
\path{GeneralizedRvachevIdentifiability.lean}, under
\path{Analysis/FabiusFunction/Lean/FabiusFunction/}.
Pinned tree \texttt{0e0368b216b302d5822890b48bd19f2337af4d04}, inspected
28 September 2026.
\url{https://github.com/VladimirReshetnikov/ProveIt/blob/0e0368b216b302d5822890b48bd19f2337af4d04/Analysis/FabiusFunction/Lean/FabiusFunction/GeneralizedRvachevIdentifiability.lean}.

\bibitem{RepoHolonomic}
\emph{Holonomic Rank, Exact Overlaps, and Non-P-Recursiveness in the
Fabius--Rvachev System}, repository research report prepared for
Vladimir Reshetnikov, dated 30 August 2026.
In \emph{ProveIt}, under
\path{Analysis/FabiusFunction/docs/semi-formalized-research-frontiers/drafts/spectra-and-arithmetic/fabius_holonomic_frontiers_report/}.
File \texttt{fabius\_holonomic\_frontiers.tex}, pinned tree
\texttt{0e0368b216b302d5822890b48bd19f2337af4d04}.
Repository: \url{https://github.com/VladimirReshetnikov/ProveIt}.
\end{thebibliography}
\end{document}
